% ============================================================================
% DescripcionDelProtocolo.tex
% ============================================================================
\documentclass[11pt,a4paper]{article}

% ============================================================================
% 1. METADATOS DEL DOCUMENTO
% ============================================================================
\newcommand{\docInstitute}{Tecnol\'ogico de Costa Rica}
\newcommand{\docSchool}{Escuela de Ingenier\'ia en Computadores}
\newcommand{\docCourse}{Seguridad de la Informaci\'on}
\newcommand{\docCode}{CE-1115}
\newcommand{\docKicker}{Taller de Criptograf\'ia Aplicada}
\newcommand{\docTitle}{SecureChat: descripci\'on del protocolo}
\newcommand{\docSubtitle}{Chat cifrado entre dos procesos locales con criptograf\'ia h\'ibrida}
\newcommand{\docTagline}{?`Qu\'e viaja por el socket y qu\'e nunca sale de cada extremo?}
\newcommand{\docKeywords}{RSA-OAEP \textbullet\ AES-256-GCM \textbullet\ Handshake \textbullet\ Nonce \textbullet\ Authentication tag}
\newcommand{\docGroup}{Trabajo en parejas}
\newcommand{\docProfessor}{MSc. Andr\'es F. Vargas Rivera}
\newcommand{\docSemester}{II Semestre, 2026}
\newcommand{\docDate}{[Fecha de entrega]}
\newcommand{\docPlace}{Cartago, Costa Rica}
\newcommand{\docShortTitle}{SecureChat --- Descripci\'on del protocolo}
\newcommand{\docRepo}{https://github.com/Adriel23456/SecureChat_Workshop}

% Integrantes: {Nombre}{Carne}
\newcommand{\docMembers}{%
  \member{Adriel S. Chaves Salazar}{2021031465}
  \member{Henry D. Nu\~nez P\'erez}{2022089224}}
\newcommand{\docAuthor}{Adriel S. Chaves Salazar, Henry D. Nu\~nez P\'erez}

% ============================================================================
% 2. IDIOMA Y FUENTES
% ============================================================================
\usepackage[T1]{fontenc}
\usepackage[utf8]{inputenc}
\usepackage{lmodern}
\usepackage[spanish,es-tabla,es-nodecimaldot,es-noquoting]{babel}
\usepackage[autostyle,spanish=mexican]{csquotes}
\usepackage{microtype}

% ============================================================================
% 3. MATEMATICA Y UNIDADES
% ============================================================================
\usepackage{amsmath}
\usepackage{amssymb}
\usepackage{amsthm}
\usepackage{mathtools}
\usepackage{siunitx}
\sisetup{per-mode=symbol,range-units=single,output-decimal-marker={,}}

% ============================================================================
% 4. COLORES
% ============================================================================
\usepackage[table]{xcolor}
\definecolor{tecnavy}{RGB}{31,56,100}
\definecolor{tecsteel}{RGB}{46,85,150}
\definecolor{teclight}{RGB}{214,227,243}
\definecolor{tecrule}{RGB}{155,184,216}
\definecolor{boxbg}{RGB}{242,246,252}
\definecolor{labelgray}{RGB}{120,120,120}
\definecolor{codebg}{RGB}{248,248,248}
\definecolor{codekw}{RGB}{0,0,180}
\definecolor{codecmt}{RGB}{0,120,60}
\definecolor{codestr}{RGB}{170,55,55}
% Clasificacion de la informacion del protocolo
\definecolor{nivelpublico}{RGB}{220,240,220}
\definecolor{nivelsecreto}{RGB}{250,215,215}

% ============================================================================
% 5. GRAFICOS Y TIKZ
% ============================================================================
\usepackage{graphicx}
\graphicspath{{img/}{figures/}{./}}
\usepackage{tikz}
\usepackage{pgfplots}
\pgfplotsset{compat=1.18}
\usetikzlibrary{arrows.meta,positioning,shapes.geometric,calc,fit,
                backgrounds,patterns,decorations.pathreplacing,babel}

% ============================================================================
% 6. TABLAS, FLOTANTES Y LISTAS
% ============================================================================
\usepackage{float}
\usepackage{array}
\usepackage{booktabs}
\usepackage{multirow}
\usepackage{makecell}
\usepackage{longtable}
\usepackage[justification=centering]{caption}
\usepackage{subcaption}
\usepackage{tabularx}
\usepackage{xltabular}
\usepackage{adjustbox}
\usepackage{enumitem}

\newcolumntype{P}[1]{>{\raggedright\arraybackslash}p{#1}}
\newcolumntype{C}[1]{>{\centering\arraybackslash}p{#1}}
\newcolumntype{R}[1]{>{\raggedleft\arraybackslash}p{#1}}
\newcolumntype{Y}{>{\raggedright\arraybackslash}X}
\renewcommand{\arraystretch}{1.15}
\setlength{\tabcolsep}{5pt}
\setlist[itemize]{leftmargin=*,noitemsep,topsep=2pt}
\setlist[enumerate]{leftmargin=*,noitemsep,topsep=2pt}

\newcommand{\contnext}[1]{%
  \multicolumn{#1}{r@{}}{\footnotesize\itshape Contin\'ua en la p\'agina siguiente}\\}
\newcommand{\contprev}[1]{%
  \multicolumn{#1}{@{}l}{\footnotesize\itshape
    Tabla \thetable{} -- viene de la p\'agina anterior}\\}

\newcommand{\thd}[1]{\cellcolor{tecnavy}\textcolor{white}{\textbf{#1}}}

\theoremstyle{definition}
\newtheorem{definicion}{Definici\'on}[section]
\theoremstyle{remark}
\newtheorem*{observacion}{Observaci\'on}

% ============================================================================
% 7. CODIGO FUENTE (LISTINGS)
% ============================================================================
\usepackage{listings}
\renewcommand{\lstlistingname}{C\'odigo}
\renewcommand{\lstlistlistingname}{\'Indice de c\'odigos}
\lstdefinestyle{code}{
  basicstyle=\ttfamily\footnotesize,
  keywordstyle=\color{codekw}\bfseries,
  commentstyle=\color{codecmt}\itshape,
  stringstyle=\color{codestr},
  backgroundcolor=\color{codebg},
  numbers=left,
  numberstyle=\tiny\color{labelgray},
  numbersep=8pt,
  frame=single,
  rulecolor=\color{tecrule},
  breaklines=true,
  breakatwhitespace=false,
  postbreak=\mbox{\textcolor{tecsteel}{$\hookrightarrow$}\space},
  showspaces=false,
  showstringspaces=false,
  tabsize=2,
  captionpos=b,
  columns=flexible,
  upquote=true,
  literate={á}{{\'a}}1 {é}{{\'e}}1 {í}{{\'i}}1 {ó}{{\'o}}1 {ú}{{\'u}}1
           {ñ}{{\~n}}1 {Á}{{\'A}}1 {É}{{\'E}}1 {Í}{{\'I}}1 {Ó}{{\'O}}1
           {Ú}{{\'U}}1 {Ñ}{{\~N}}1 {¿}{{?`}}1 {¡}{{!`}}1
}
\lstset{style=code}

% ============================================================================
% 8. DISENO DE PAGINA, TITULOS Y ENCABEZADOS
% ============================================================================
\usepackage[a4paper,margin=1in,headheight=15pt]{geometry}
\usepackage{setspace}
\setstretch{1.15}
\usepackage{ragged2e}

\usepackage{titlesec}
\titleformat{\section}
  {\Large\bfseries\color{tecnavy}}{\thesection.}{0.6em}{}
\titleformat{\subsection}
  {\large\bfseries\color{tecsteel}}{\thesubsection.}{0.6em}{}
\titleformat{\subsubsection}
  {\normalsize\bfseries\color{tecsteel}}{\thesubsubsection.}{0.6em}{}

\usepackage{tocloft}
\renewcommand{\cfttoctitlefont}{\hfill\Large\bfseries\color{tecnavy}}
\renewcommand{\cftaftertoctitle}{\hfill}

\usepackage{fancyhdr}
\usepackage{lastpage}
\pagestyle{fancy}
\fancyhf{}
\fancyhead[L]{\footnotesize\textsc{\docCode}}
\fancyhead[R]{\footnotesize\docShortTitle}
\fancyfoot[C]{\footnotesize P\'agina \thepage\ de \pageref{LastPage}}
\renewcommand{\headrulewidth}{0.4pt}
\renewcommand{\footrulewidth}{0pt}

\counterwithin{figure}{section}
\counterwithin{table}{section}
\AtBeginDocument{\counterwithin{lstlisting}{section}}

% ============================================================================
% 9. CITAS (APA) Y ENLACES
% ============================================================================
\usepackage[authoryear,round]{natbib}
\renewcommand{\bibsection}{}
\usepackage{hyperref}
\hypersetup{
  colorlinks=true,
  linkcolor=tecnavy,
  citecolor=tecsteel,
  urlcolor=tecsteel,
  bookmarksnumbered=true,
  pdftitle={\docTitle},
  pdfauthor={\docAuthor},
  pdfsubject={\docCourse\ (\docCode) -- \docKicker},
  pdflang={es}
}
\usepackage[spanish]{cleveref}

% ============================================================================
% 10. COMANDOS PROPIOS
% ============================================================================
\newcommand{\missingfilebox}[1]{%
  \par\medskip\noindent
  \fcolorbox{red!60!black}{red!5}{%
    \parbox{\dimexpr\linewidth-2\fboxsep-2\fboxrule}{%
      \centering\ttfamily\small ARCHIVO FALTANTE: \detokenize{#1}}}%
  \par\medskip}

\newcommand{\missingimagebox}[1]{%
  \begingroup
  \setlength{\fboxsep}{0pt}%
  \fcolorbox{tecrule}{boxbg}{%
    \parbox[c][2.8cm][c]{6cm}{\centering\ttfamily\scriptsize
      IMAGEN NO ENCONTRADA\\[3pt]\detokenize{#1}}}%
  \endgroup}

\newcommand{\safeimage}[2][width=0.75\linewidth]{%
  \IfFileExists{#2}{\includegraphics[#1]{#2}}{%
    \IfFileExists{img/#2}{\includegraphics[#1]{img/#2}}{%
      \IfFileExists{figures/#2}{\includegraphics[#1]{figures/#2}}{%
        \PackageWarning{tecdoc}{Falta la imagen #2}%
        \missingimagebox{#2}}}}}

\newcommand{\safelisting}[2][]{%
  \IfFileExists{#2}{\lstinputlisting[#1]{#2}}{%
    \PackageWarning{tecdoc}{Falta el codigo #2}%
    \missingfilebox{#2}}}

\newcommand{\notebox}[2][Nota]{%
  \par\medskip\noindent
  \fcolorbox{tecrule}{boxbg}{%
    \parbox{\dimexpr\linewidth-2\fboxsep-2\fboxrule}{%
      \vspace{2pt}\textbf{\color{tecnavy}#1.}\ #2\vspace{2pt}}}%
  \par\medskip}

\newcommand{\code}[1]{\texttt{#1}}

\newcommand{\pendiente}[1]{%
  \textcolor{codestr}{\textbf{[PENDIENTE:} \textit{#1}\textbf{]}}}
\newcommand{\pend}{\textcolor{codestr}{\textbf{[PEND.]}}}

% Etiquetas de clasificacion
\newcommand{\nivel}[2]{\colorbox{#1}{\strut\footnotesize\textbf{#2}}}
\newcommand{\publico}{\nivel{nivelpublico}{P\'ublico}}
\newcommand{\secreto}{\nivel{nivelsecreto}{Secreto}}

% ============================================================================
% 11. PORTADA
% ============================================================================
\newcommand{\member}[2]{%
  {\small\bfseries #1} & {\footnotesize\color{labelgray} #2}\\[3pt]}

\newcommand{\makecover}{%
  \thispagestyle{empty}
  \begingroup
  \setlength{\parindent}{0pt}
  \setlength{\parskip}{0pt}

  \begin{tikzpicture}[remember picture, overlay]
    \fill[tecnavy]
      (current page.north west) rectangle
      ([yshift=-5.6cm] current page.north east);
    \fill[tecsteel]
      ([yshift=-5.6cm] current page.north west) rectangle
      ([yshift=-5.8cm] current page.north east);
    \node[anchor=north, inner sep=0pt]
      at ([yshift=-0.7cm] current page.north)
      {\safeimage[width=1.1in]{TUlogo.png}};
    \node[anchor=center]
      at ([yshift=-4.3cm] current page.north)
      {\parbox{0.80\textwidth}{\centering
        {\large\bfseries\color{white} \docInstitute}\\[4pt]
        {\small\color{white!80!tecnavy} \docSchool}\\[3pt]
        {\footnotesize\color{white!65!tecnavy}
          \docCourse\ \textbullet\ \docCode}}};
  \end{tikzpicture}

  \vspace*{2.4cm}
  \begin{center}
    {\large\bfseries \docKicker}
  \end{center}

  \vspace{0.1cm}
  \begin{center}
    \colorbox{teclight}{%
      \parbox{0.75\textwidth}{\centering\vspace{5pt}
        {\footnotesize\color{tecsteel} \docSubtitle}%
      \vspace{5pt}}}\\[0.6cm]
    {\LARGE\bfseries\color{tecnavy} \docTitle}\\[0.30cm]
    {\large\itshape \docTagline}\\[0.45cm]
    {\color{tecrule}\rule{0.5\textwidth}{0.4pt}}\\[0.45cm]
    {\normalsize \docKeywords}\\[0.35cm]
    {\footnotesize\color{labelgray} Repositorio: \url{\docRepo}}
  \end{center}

  \vfill

  \noindent
  \fcolorbox{tecrule}{boxbg}{%
    \parbox{\dimexpr\textwidth-2\fboxsep-2\fboxrule}{%
      \vspace{0.25cm}
      \centering
      {\footnotesize\color{labelgray} PROFESOR}\\[0.2cm]
      {\small\bfseries \docProfessor}
      \vspace{0.25cm}}}

  \vspace{0.5cm}

  \noindent
  \fcolorbox{tecrule}{boxbg}{%
    \parbox{\dimexpr\textwidth-2\fboxsep-2\fboxrule}{%
      \vspace{0.25cm}
      \centering
      {\footnotesize\color{labelgray} INTEGRANTES --- \MakeUppercase{\docGroup}}\\[0.2cm]
      \begin{tabular*}{0.88\textwidth}{@{}l@{\extracolsep{\fill}}r@{}}
        {\footnotesize\color{labelgray} Nombre} &
        {\footnotesize\color{labelgray} Carn\'e}\\[2pt]
        \docMembers
      \end{tabular*}
      \vspace{0.15cm}}}

  \vspace{0.4cm}
  {\color{tecsteel}\noindent\rule{\textwidth}{0.8pt}}

  \vspace{-0.5cm}
  \begin{center}
    {\footnotesize\color{labelgray}
      \docSemester\ \textbullet\ \docDate\ \textbullet\ \docPlace}
  \end{center}

  \endgroup
  \newpage}

% ============================================================================
% 12. CUERPO DEL DOCUMENTO
% ============================================================================
\begin{document}

\pagenumbering{roman}
\makecover

\tableofcontents
\newpage

\pagenumbering{arabic}
\justifying

%#############################################################################
% SECCION 1 -- INTRODUCCION Y OBJETIVO
%#############################################################################
\section{Introducción y objetivo}
\label{sec:intro}

SecureChat es una aplicación de consola que mantiene una conversación
bidireccional entre dos procesos que se ejecutan en la misma computadora,
conectados por un socket TCP sobre \code{localhost}. La criptografía
\textbf{asimétrica} se usa para establecer la sesión y la criptografía
\textbf{simétrica autenticada} protege cada mensaje.

\pendiente{Resumen de la solución final en un párrafo: lenguaje, biblioteca,
algoritmos, tamaños de llave y resultado de la demostración.}

\subsection{Alcance}
\label{subsec:alcance}

\pendiente{Qué cubre y qué no cubre la solución (por ejemplo, autenticación
de la llave pública, persistencia, más de dos participantes).}

\subsection{Entorno de desarrollo}
\label{subsec:entorno}

\begin{table}[H]
  \centering
  \caption{Entorno de desarrollo y demostración}
  \label{tab:entorno}
  \small
  \begin{tabularx}{\linewidth}{@{}P{4cm} Y@{}}
    \toprule
    \textbf{Elemento} & \textbf{Detalle} \\
    \midrule
    Sistema operativo & Windows \pend \\
    Lenguaje          & Python \pend \\
    Biblioteca criptográfica & \code{cryptography} \pend \\
    Comunicación      & Sockets TCP/IP sobre \code{127.0.0.1}, puerto \pend \\
    Repositorio       & \url{\docRepo} \\
    \bottomrule
  \end{tabularx}
\end{table}
% ---------- FIN SECCION 1 ----------

%#############################################################################
% SECCION 2 -- ARQUITECTURA GENERAL
%#############################################################################
\section{Arquitectura general}
\label{sec:arquitectura}

El mismo programa se ejecuta dos veces, una instancia por participante
(\cref{fig:arquitectura}). Bob espera conexiones y Alice inicia la conexión.

\begin{figure}[H]
  \centering
  \begin{tikzpicture}[font=\small,
      inst/.style={draw=tecnavy, fill=boxbg, rounded corners=3pt,
                   text width=4.3cm, minimum height=1.3cm, align=center},
      mat/.style={draw=tecsteel, fill=teclight, rounded corners=3pt,
                  text width=4.6cm, minimum height=1.0cm, align=center,
                  font=\footnotesize}]
    \node[inst] (a) at (0,0) {\textbf{Terminal A -- Alice}\\Instancia de SecureChat\\
      {\scriptsize (inicia la conexión)}};
    \node[inst] (b) at (8.6,0) {\textbf{Terminal B -- Bob}\\Instancia de SecureChat\\
      {\scriptsize (espera conexiones)}};
    \draw[{Stealth[length=2mm]}-{Stealth[length=2mm]}, thick, tecsteel] (a) --
      node[above, font=\footnotesize]{socket TCP}
      node[below, font=\footnotesize]{\code{127.0.0.1}} (b);
    \node[mat, below=1.0cm of a] (ma) {Material criptográfico local\\
      llave privada, llave pública y clave de sesión};
    \node[mat, below=1.0cm of b] (mb) {Material criptográfico local\\
      llave privada, llave pública y clave de sesión};
    \draw[-{Stealth[length=2mm]}, thick, tecnavy] (a) -- (ma);
    \draw[-{Stealth[length=2mm]}, thick, tecnavy] (b) -- (mb);
  \end{tikzpicture}
  \caption{Arquitectura de SecureChat: dos instancias del mismo programa
  conectadas por TCP en la misma computadora.}
  \label{fig:arquitectura}
\end{figure}

\subsection{Organización del código}
\label{subsec:codigo}

\pendiente{Módulos finales y responsabilidad de cada uno (red, framing,
llaves, handshake, cifrado, interfaz de consola).}

\subsection{Concurrencia}
\label{subsec:concurrencia}

\pendiente{Cómo se reciben mensajes mientras el usuario escribe (hilo
receptor y entrada de teclado) y por qué en Windows no se usa
\code{select()} sobre \code{stdin}.}
% ---------- FIN SECCION 2 ----------

%#############################################################################
% SECCION 3 -- MATERIAL CRIPTOGRAFICO
%#############################################################################
\section{Material criptográfico}
\label{sec:material}

Cada participante dispone de su propio par de llaves RSA, generado y
guardado dentro de su carpeta (\code{srcAlice/keys/} y \code{srcBob/keys/}).
Ninguna de estas carpetas se versiona: la regla \code{keys/} del
\code{.gitignore} las excluye del repositorio.

\subsection{Par de llaves asimétricas}
\label{subsec:rsa}

\begin{table}[H]
  \centering
  \caption{Parámetros del par de llaves asimétricas}
  \label{tab:rsa}
  \small
  \begin{tabularx}{\linewidth}{@{}P{4.6cm} Y@{}}
    \toprule
    \textbf{Parámetro} & \textbf{Valor} \\
    \midrule
    Algoritmo                   & RSA \\
    Tamaño del módulo           & 3072 bits (se rechaza al cargar cualquier llave menor a 2048 bits) \\
    Exponente público           & $e = 65537$ \\
    Generación                  & \code{rsa.generate\_private\_key} de \code{cryptography}, con el generador aleatorio seguro del sistema operativo \\
    Relleno para cifrar         & OAEP con SHA-256 y MGF1(SHA-256) (uso en el handshake, \cref{sec:handshake}) \\
    Relleno para firmar         & PSS con SHA-256 y MGF1(SHA-256) \\
    Formato de la llave privada & PEM, PKCS\#8, \textbf{cifrada} con una frase de paso (\code{BestAvailableEncryption}) \\
    Formato de la llave pública & PEM, \emph{SubjectPublicKeyInfo} (X.509) \\
    Frase de paso               & Mínimo 8 caracteres; se pide sin eco en pantalla (\code{getpass}) y nunca se guarda \\
    Archivos                    & \code{keys/private\_key.pem} y \code{keys/public\_key.pem} dentro de la carpeta de cada participante \\
    \bottomrule
  \end{tabularx}
\end{table}

Se eligió RSA-3072 en lugar del mínimo de 2048 bits porque ofrece alrededor
de 128 bits de seguridad, un nivel que se mantiene vigente a largo plazo.
El costo es una generación de llaves algo más lenta (uno o dos segundos),
que ocurre una sola vez. El exponente $e = 65537$ es el valor estándar:
verificar y cifrar con él es rápido y no presenta las debilidades de los
exponentes muy pequeños.

\subsection{Llave pública y llave privada: qué hace cada una}
\label{subsec:publicaprivada}

Las dos llaves están relacionadas matemáticamente: lo que una hace, solo la
otra lo deshace. La llave pública $(n, e)$ se deriva de la privada, pero el
camino inverso exige factorizar $n$, lo cual es inviable para un módulo de
3072 bits. Por eso:

\begin{itemize}
  \item \textbf{La llave pública puede compartirse.} Sirve para
        \emph{cifrar hacia} su dueño y para \emph{verificar} sus firmas.
        Conocerla no permite descifrar ni firmar.
  \item \textbf{La llave privada debe permanecer bajo control exclusivo de
        su dueño.} Es la única que puede \emph{descifrar} lo cifrado con la
        pública y \emph{firmar}. Nunca se muestra en pantalla, nunca se
        escribe sin cifrar y nunca viaja por el socket.
  \item \textbf{La llave pública no reemplaza a la privada.} Con la pública
        de Bob, Alice puede proteger la clave de sesión para que solo Bob la
        lea, pero no puede leerla ella misma una vez cifrada ni hacerse pasar
        por Bob. Si la pública bastara para descifrar, cualquiera que la
        conociera (es decir, todos) podría hacerlo.
\end{itemize}

\notebox[Prueba de correspondencia]{La opción \emph{Verificar mis llaves}
del menú carga ambos archivos y ejecuta tres comprobaciones: (1) el módulo
$n$ y el exponente $e$ derivados de la privada coinciden con los de la
pública; (2) una firma PSS hecha con la privada es aceptada por la pública;
y (3) un valor aleatorio cifrado con OAEP usando la pública se recupera con
la privada. El valor de prueba es aleatorio, se descarta y solo se informa
\enquote{aprobado} o \enquote{fallido}.}

\subsection{Huella de la llave pública}
\label{subsec:huella}

La huella es el resumen SHA-256 de la codificación DER de la llave pública:
32 bytes que se muestran como \code{SHA256:ab:cd:...}. Una llave de 3072
bits ocupa cientos de bytes y no se puede comparar a simple vista; la
huella sí. Además, si un atacante sustituye la llave pública durante el
handshake, la huella cambia por completo.

La huella es información pública: mostrarla no revela ningún secreto. Es lo
único de las llaves que el programa muestra en pantalla.

\subsection{Distribución de las llaves públicas}
\label{subsec:distribucion}

\notebox[Decisión]{Se adoptó un modelo mixto. La llave pública viaja por
el socket durante el handshake, pero antes de usarla el receptor compara su
huella SHA-256 con una huella de confianza registrada previamente. Para
registrarla, cada participante muestra su huella en su terminal y el otro
la copia por un canal independiente de la red (\emph{fuera de banda}). Si
las huellas no coinciden, el handshake se aborta. Así se cierra la puerta
al atacante en el medio que sustituye la llave pública, un riesgo que la
criptografía por sí sola no resuelve si no se autentica la llave.}

Las huellas de confianza se guardan en
\code{srcBob/keys/trusted\_alice\_fingerprint.txt} y
\code{srcAlice/keys/trusted\_bob\_fingerprint.txt}.

\subsection{Uso del par de llaves de cada participante}
\label{subsec:usopar}

\begin{itemize}
  \item \textbf{Par de Bob:} Alice cifra la clave de sesión con la llave
        pública de Bob (RSA-OAEP) y solo Bob puede recuperarla con su
        privada. Que Bob logre confirmar la sesión con esa clave demuestra
        que posee la privada correspondiente a la huella registrada.
  \item \textbf{Par de Alice:} Alice firma su mensaje del handshake con su
        llave privada (RSA-PSS) y Bob verifica la firma con la llave pública
        de Alice, previamente autenticada por su huella. Así Bob sabe que la
        clave de sesión la generó Alice y no un tercero.
\end{itemize}

Con este esquema la autenticación es \textbf{mutua} y cada par de llaves
cumple una función concreta. \pendiente{Confirmar con la implementación
del Paso 3.}

\subsection{Clave de sesión}
\label{subsec:clavesesion}

\pendiente{Cómo y cuándo se genera la clave de sesión (256 bits, generador
aleatorio criptográfico, una nueva por sesión, nunca fija en el código) y
cuánto vive. Se completa en el Paso 3.}

\subsection{Qué es público, qué es secreto y qué viaja por la red}
\label{subsec:clasificacion}

\begin{table}[H]
  \centering
  \caption{Clasificación del material y de los datos del protocolo}
  \label{tab:clasificacion}
  \small
  \begin{tabularx}{\linewidth}{@{}P{3.6cm} Y P{2.7cm} l@{}}
    \toprule
    \textbf{Elemento} & \textbf{Dónde vive} & \textbf{¿Viaja por el socket?} & \textbf{Nivel} \\
    \midrule
    Llave privada (Alice, Bob) & \code{keys/private\_key.pem}, cifrada; en memoria solo mientras se usa & Nunca & \secreto \\
    Frase de paso              & Solo en la memoria de su dueño & Nunca & \secreto \\
    Llave pública              & \code{keys/public\_key.pem} & Sí, en el handshake & \publico \\
    Huella de la llave pública & Pantalla y archivo de confianza del par & No (se compara fuera de banda) & \publico \\
    Clave de sesión            & Memoria de ambos procesos & Solo cifrada con RSA-OAEP & \secreto \\
    Plaintext del usuario      & Solo en la terminal local & Nunca (tras el handshake) & \secreto \\
    Nonce                      & Generado por mensaje & Sí, junto al ciphertext & \publico \\
    Ciphertext                 & En tránsito & Sí & \publico \\
    Authentication tag         & En tránsito & Sí & \publico \\
    Metadatos (puertos, tamaños, horas) & Encabezados TCP/IP & Sí & \publico \\
    \bottomrule
  \end{tabularx}
\end{table}
% ---------- FIN SECCION 3 ----------

%#############################################################################
% SECCION 4 -- HANDSHAKE
%#############################################################################
\section{Establecimiento de la sesión (handshake)}
\label{sec:handshake}

El handshake se completa antes de permitir cualquier mensaje de usuario
(\cref{fig:handshake}).

\begin{figure}[H]
  \centering
  \begin{tikzpicture}[font=\footnotesize,
      msg/.style={-{Stealth[length=2mm]}, thick, tecsteel},
      nota/.style={draw=tecrule, fill=boxbg, rounded corners=2pt,
                   align=center, font=\scriptsize, inner sep=3pt}]
    \node[font=\bfseries, text=tecnavy] at (0,0) {Alice};
    \node[font=\bfseries, text=labelgray] at (6,0) {Socket / red local};
    \node[font=\bfseries, text=tecnavy] at (12,0) {Bob};
    \foreach \x in {0,6,12}
      \draw[dashed, labelgray] (\x,-0.3) -- (\x,-7.4);
    \draw[msg] (0,-1.0) -- node[above]{1. Conexión TCP / saludo inicial \pend} (12,-1.0);
    \draw[msg] (12,-2.0) -- node[above]{2. Llave pública de Bob \pend} (0,-2.0);
    \node[nota] at (0,-2.9) {Genera $K_s$ (256 bits)\\con CSPRNG};
    \draw[msg] (0,-3.9) -- node[above]{3. $\mathrm{RSA\text{-}OAEP}_{\mathit{pub}_B}(K_s)$} (12,-3.9);
    \node[nota] at (12,-4.8) {Descifra con $\mathit{priv}_B$\\y recupera $K_s$};
    \draw[msg, codecmt] (12,-5.7) -- node[above]{4. Confirmación de sesión \pend} (0,-5.7);
    \node[draw=codecmt, fill=codecmt!8, rounded corners=2pt, align=center,
          text width=8cm, inner sep=4pt] at (6,-6.8)
      {\textbf{Ambos poseen $K_s$.} A partir de aquí todo mensaje se protege
       con AES-256-GCM.};
  \end{tikzpicture}
  \caption{Secuencia del handshake (versión preliminar; se actualizará con
  los mensajes reales).}
  \label{fig:handshake}
\end{figure}

\subsection{Mensajes del handshake}
\label{subsec:mensajeshs}

\begin{table}[H]
  \centering
  \caption{Mensajes del handshake}
  \label{tab:mensajeshs}
  \small
  \begin{tabularx}{\linewidth}{@{}c l l Y@{}}
    \toprule
    \textbf{\#} & \textbf{Emisor} & \textbf{Tipo} & \textbf{Contenido y protección} \\
    \midrule
    1 & Alice & \pend & \pend \\
    2 & Bob   & \pend & \pend \\
    3 & Alice & \pend & \pend \\
    4 & Bob   & \pend & \pend \\
    \bottomrule
  \end{tabularx}
\end{table}

\subsection{Confirmación de la sesión}
\label{subsec:confirmacion}

\pendiente{Cómo se prueba que ambos extremos tienen la misma clave de
sesión y qué se muestra en pantalla al establecerse la sesión segura.}

\subsection{Manejo de errores en el handshake}
\label{subsec:erroreshs}

\pendiente{Qué ocurre si la llave pública es inválida, el descifrado falla,
la confirmación no coincide o se agota el tiempo. Nunca se continúa sin
cifrado.}

\notebox[Decisión pendiente]{Autenticación de la llave pública de Bob.
Sin verificarla, un atacante activo en el medio podría sustituirla.
\pendiente{Decisión adoptada (por ejemplo, verificación de la huella
SHA-256 fuera de banda) y justificación.}}
% ---------- FIN SECCION 4 ----------

%#############################################################################
% SECCION 5 -- PROTECCION DE MENSAJES
%#############################################################################
\section{Protección de los mensajes}
\label{sec:mensajes}

\subsection{AES-256-GCM}
\label{subsec:gcm}

\begin{table}[H]
  \centering
  \caption{Parámetros del cifrado simétrico autenticado}
  \label{tab:gcm}
  \small
  \begin{tabularx}{\linewidth}{@{}P{4.2cm} Y@{}}
    \toprule
    \textbf{Parámetro} & \textbf{Valor} \\
    \midrule
    Algoritmo           & AES-256-GCM (cifrado autenticado, AEAD) \\
    Tamaño de clave     & 256 bits \\
    Tamaño del nonce    & 96 bits \\
    Tamaño del tag      & 128 bits \\
    Datos asociados (AAD) & \pend \\
    Codificación del texto & \pend \\
    \bottomrule
  \end{tabularx}
\end{table}

\subsection{Generación del nonce}
\label{subsec:nonce}

\pendiente{Estrategia adoptada (contador por dirección o aleatorio), cómo
se evita la repetición con la misma clave considerando que ambos sentidos
comparten la clave de sesión, y qué pasa si el contador se agota.}

\subsection{Verificación en el receptor}
\label{subsec:verificacion}

\pendiente{El receptor verifica el tag antes de aceptar el plaintext; si
falla, el mensaje se rechaza y no se muestra contenido.}

\subsection{Protección contra repetición}
\label{subsec:replay}

\pendiente{Uso de números de secuencia autenticados o justificación si no
se implementa.}
% ---------- FIN SECCION 5 ----------

%#############################################################################
% SECCION 6 -- FORMATO DE LOS MENSAJES TRANSMITIDOS
%#############################################################################
\section{Formato de los mensajes transmitidos}
\label{sec:formato}

TCP entrega un flujo de bytes sin límites de mensaje, por lo que cada
mensaje se delimita con un prefijo de longitud. \pendiente{Confirmar.}

\begin{table}[H]
  \centering
  \caption{Estructura de una trama cifrada (propuesta preliminar)}
  \label{tab:trama}
  \small
  \begin{tabularx}{\linewidth}{@{}l l Y@{}}
    \toprule
    \textbf{Campo} & \textbf{Tamaño} & \textbf{Descripción} \\
    \midrule
    Longitud     & \pend & Tamaño del resto de la trama \\
    Tipo         & \pend & Tipo de mensaje \\
    Secuencia    & \pend & Número de mensaje por dirección \\
    Nonce        & 12 bytes & Único por cifrado con la misma clave \\
    Ciphertext   & variable & Mensaje cifrado \\
    Tag          & 16 bytes & Authentication tag de GCM \\
    \bottomrule
  \end{tabularx}
\end{table}

\pendiente{Diagrama de la trama en bytes y ejemplo real capturado.}

\notebox[Codificación no es cifrado]{La representación en hexadecimal o
Base64 se usa solo para mostrar o transportar bytes. No ofrece
confidencialidad.}
% ---------- FIN SECCION 6 ----------

%#############################################################################
% SECCION 7 -- CONVERSACION Y CIERRE
%#############################################################################
\section{Conversación bidireccional y cierre de sesión}
\label{sec:conversacion}

\subsection{Flujo de un mensaje}
\label{subsec:flujo}

Mensaje legible $\rightarrow$ AES-GCM cifra y autentica $\rightarrow$
ciphertext en tránsito $\rightarrow$ verificación y descifrado
$\rightarrow$ mensaje recuperado.

\pendiente{Diagrama de secuencia de un mensaje en cada sentido.}

\subsection{Presentación en consola}
\label{subsec:consola}

\begin{table}[H]
  \centering
  \caption{Distinción visual en la consola}
  \label{tab:consola}
  \small
  \begin{tabularx}{\linewidth}{@{}l Y@{}}
    \toprule
    \textbf{Etiqueta} & \textbf{Significado} \\
    \midrule
    \pend & Mensaje local (plaintext antes de cifrar) \\
    \pend & Datos cifrados enviados \\
    \pend & Datos cifrados recibidos \\
    \pend & Mensaje recuperado tras verificar el tag \\
    \pend & Error o mensaje rechazado \\
    \bottomrule
  \end{tabularx}
\end{table}

\subsection{Cierre controlado}
\label{subsec:cierre}

\pendiente{Comando de salida, mensaje de cierre cifrado y comportamiento
del otro extremo.}
% ---------- FIN SECCION 7 ----------

%#############################################################################
% SECCION 8 -- DETECCION DE MANIPULACION
%#############################################################################
\section{Detección de manipulación}
\label{sec:manipulacion}

\pendiente{Mecanismo de la prueba controlada (alteración en el emisor o
proxy intermedio), qué byte se altera y por qué.}

\pendiente{Resultado: el receptor detecta el tag inválido, rechaza el
mensaje y no muestra plaintext. Política tras el rechazo (descartar o
cerrar la sesión).}

\begin{figure}[H]
  \centering
  \safeimage[width=0.9\linewidth]{05_manipulacion.png}
  \caption{Rechazo de un ciphertext alterado.}
  \label{fig:manipulacion}
\end{figure}
% ---------- FIN SECCION 8 ----------

%#############################################################################
% SECCION 9 -- ANALISIS DE SEGURIDAD
%#############################################################################
\section{Análisis de seguridad}
\label{sec:analisis}

\subsection{Modelo de adversario}
\label{subsec:adversario}

\pendiente{Adversario pasivo (escucha el socket) y activo (modifica,
repite o inyecta), en la línea del modelo Dolev-Yao visto en clase.}

\subsection{Cumplimiento de los requisitos mínimos}
\label{subsec:cumplimiento}

{\small
\begin{xltabular}{\linewidth}{@{}c P{5.2cm} Y P{2.6cm}@{}}
  \caption{Requisitos mínimos del taller y su cumplimiento}
  \label{tab:cumplimiento}\\
  \toprule
  & \textbf{Requisito} & \textbf{Cómo se cumple} & \textbf{Evidencia} \\
  \midrule
  \endfirsthead
  \contprev{4}
  \toprule
  & \textbf{Requisito} & \textbf{Cómo se cumple} & \textbf{Evidencia} \\
  \midrule
  \endhead
  \midrule
  \contnext{4}
  \endfoot
  \bottomrule
  \endlastfoot
  $\checkmark$ & Python 3 y sockets TCP/IP & \pend & \pend \\
  $\checkmark$ & Biblioteca criptográfica reconocida & \pend & \pend \\
  $\checkmark$ & Handshake basado en criptografía asimétrica & \pend & \pend \\
  $\checkmark$ & Clave de sesión dinámica, no fija en el código & \pend & \pend \\
  $\checkmark$ & Ninguna llave privada viaja por el socket & \pend & \pend \\
  $\checkmark$ & Mensajes cifrados con esquema autenticado & \pend & \pend \\
  $\checkmark$ & Nonce correcto y no repetido & \pend & \pend \\
  $\checkmark$ & Rechazo de tags inválidos & \pend & \pend \\
  $\checkmark$ & Evidencia de plaintext, ciphertext enviado y recibido, y plaintext recuperado & \pend & \pend \\
  $\times$ & Base64 o hexadecimal no sustituyen al cifrado & \pend & \pend \\
  $\times$ & No se usa ECB & \pend & \pend \\
\end{xltabular}}

\subsection{Limitaciones conocidas}
\label{subsec:limitaciones}

\pendiente{Por ejemplo: autenticación de la llave pública, ausencia de
secreto hacia adelante (forward secrecy) al transportar la clave con RSA,
análisis de tráfico (tamaños y tiempos visibles), y la imposibilidad de
borrar con certeza secretos de la memoria en Python.}
% ---------- FIN SECCION 9 ----------

%#############################################################################
% SECCION 10 -- PREGUNTAS DE CIERRE
%#############################################################################
\section{Preguntas de cierre}
\label{sec:preguntas}

\begin{enumerate}[itemsep=6pt]
  \item \textbf{¿Qué problema resuelve la criptografía asimétrica durante el handshake?}\\
        \pendiente{Respuesta.}
  \item \textbf{¿Por qué la conversación utiliza criptografía simétrica después del handshake?}\\
        \pendiente{Respuesta.}
  \item \textbf{¿Qué información puede observar un tercero que capture los paquetes de la conversación?}\\
        \pendiente{Respuesta.}
  \item \textbf{¿Por qué el nonce puede viajar junto al ciphertext y por qué no debe repetirse?}\\
        \pendiente{Respuesta.}
  \item \textbf{¿Qué función cumple el authentication tag?}\\
        \pendiente{Respuesta.}
  \item \textbf{¿Qué ocurriría si la clave de sesión se transmitiera sin protección?}\\
        \pendiente{Respuesta.}
  \item \textbf{¿Qué ocurriría si un atacante obtiene una de las claves privadas?}\\
        \pendiente{Respuesta.}
\end{enumerate}
% ---------- FIN SECCION 10 ----------

%#############################################################################
% SECCION 11 -- EVIDENCIAS DE LA DEMOSTRACION
%#############################################################################
\section{Evidencias de la demostración}
\label{sec:evidencias}

\begin{figure}[H]
  \centering
  \safeimage[width=0.9\linewidth]{01_conexion_handshake.png}
  \caption{Dos instancias en ejecución y handshake completado.}
  \label{fig:ev-handshake}
\end{figure}

\begin{figure}[H]
  \centering
  \safeimage[width=0.9\linewidth]{02_conversacion.png}
  \caption{Conversación con al menos tres mensajes por sentido.}
  \label{fig:ev-conversacion}
\end{figure}

\begin{figure}[H]
  \centering
  \safeimage[width=0.9\linewidth]{03_plaintext_vs_ciphertext.png}
  \caption{Comparación entre el mensaje legible y el ciphertext transmitido.}
  \label{fig:ev-ciphertext}
\end{figure}

\begin{figure}[H]
  \centering
  \safeimage[width=0.9\linewidth]{04_captura_red.png}
  \caption{(Opcional) Captura del tráfico de \code{localhost}.}
  \label{fig:ev-red}
\end{figure}

\pendiente{Enlace al video de demostración.}
% ---------- FIN SECCION 11 ----------

% ---------- Referencias (APA) ----------
\newpage
\section*{Referencias}
\addcontentsline{toc}{section}{Referencias}
\IfFileExists{referencias.bib}{%
  \bibliographystyle{apalike}
  \bibliography{referencias}}{\missingfilebox{referencias.bib}}

% ---------- Anexos ----------
\newpage
\appendix
\renewcommand{\thesection}{\Alph{section}}

%#############################################################################
% ANEXO A -- EJECUCION Y ESTRUCTURA DEL REPOSITORIO
%#############################################################################
\section{Ejecución y estructura del repositorio}
\label{anx:ejecucion}

\pendiente{Estructura final del repositorio y comandos para instalar,
generar llaves e iniciar cada instancia en Windows (resumen del README).}
% ---------- FIN ANEXO A ----------

%#############################################################################
% ANEXO B -- DECLARACION DE USO DE IA
%#############################################################################
\section{Declaración de uso de IA}
\label{anx:ia}

\pendiente{Herramientas usadas, para qué y cómo se verificó el resultado.}
% ---------- FIN ANEXO B ----------

\end{document}